\SourceLeaf{71}%
% A544-P0071-D02
6.\quad Numerus $5$ est non-residuum quadraticum numeri $103$ (theor. fund. et 5).

% A544-P0071-D03
7.\quad Numerus $2$ est non-residuum quadraticum numeri $5$ (theor. III).

% A544-P0071-D04
8.\quad Numerus $7$ est non-residuum quadraticum numeri $5$ (I et 7).

% A544-P0071-D05
9.\quad Numerus $5$ est non-residuum quadraticum numeri $7$ (theor. fund. et 8).

% A544-P0071-D06
10.\quad Numerus $103$ est non-residuum quadraticum numeri $7$ (I et 9).

% A544-P0071-D07
11.\quad Numerus $7$ est residuum quadraticum numeri $103$ (theor. fund. et 10).

% A544-P0071-D08
12.\quad Numerus $70$ est non-residuum quadraticum numeri $103$ (II, 1, 6, 11).

% A544-P0071-D09
13.\quad Numerus $379$ est non-residuum quadraticum numeri $103$ (I et 12).

% A544-P0071-D10
14.\quad Numerus $103$ est residuum quadraticum numeri $379$ (theor. fund. et 13).

% A544-P0071-D11
In sequentibus brevitatis caussa utemur signo in \textit{Comment. Gotting. Vol. XVI} introducto. Scilicet per $[x]$ denotabimus quantitatem $x$ ipsam, quoties $x$ est integer, sive integrum proxime minorem quam $x$, quoties $x$ est quantitas fracta, ita ut $x-[x]$ semper fiat quantitas non negativa unitate minor.

% A544-P0071-D12
\Article{3}

% A544-P0071-D13
\textsc{Problema.}\quad \textit{Denotantibus $a,\ b$ integros positivos inter se primos, et posito $[\frac12a]=a'$, invenire aggregatum}

% A544-P0071-D14
\[
\gbr{\frac ba}+\gbr{\frac{2b}a}+\gbr{\frac{3b}a}+\gbr{\frac{4b}a}+\text{etc.}+\gbr{\frac{a'b}a}
\]

% A544-P0071-D15
\textsc{Sol.}\quad Designemus brevitatis caussa huiusmodi aggregatum per $\varphi(a,b)$, ita ut etiam fiat

% A544-P0071-D16
\[
\varphi(b,a)=\gbr{\frac ab}+\gbr{\frac{2a}b}+\gbr{\frac{3a}b}+\text{etc.}+\gbr{\frac{b'a}b}
\]

% A544-P0071-D17
si statuimus $[\frac12b]=b'$. In demonstratione tertia theorematis fundamentalis ostensum est, pro casu eo, ubi $a$ et $b$ sunt impares, fieri

% A544-P0071-D18
\[
\varphi(a,b)+\varphi(b,a)=a'b'
\]

% A544-P0071-D19
facileque eandem methodum sequendo veritas huius propositionis ad eum quoque casum extenditur, ubi alteruter numerorum $a,\ b$ est impar, uti illic iam addigitavimus. Dividatur, ad instar methodi, per quam duorum integrorum divisor communis maximus investigatur, $a$ per $b$, sitque $\beta$ quotiens atque $c$ residuum; dein dividatur $b$ per $c$ et sic porro, ita ut habeantur aequationes

\SourceLeaf{72}%
% A544-P0072-D02
\[
\begin{aligned}
a&=\beta b+c\\
b&=\gamma c+d\\
c&=\delta d+e\\
d&=\varepsilon e+f\ \text{etc.}
\end{aligned}
\]

% A544-P0072-D03
Hoc modo in serie numerorum continuo decrescentium $b,c,d,e,f$ etc. tandem ad unitatem perveniemus, quum per hyp. $a$ et $b$ sint inter se primi, ita ut aequatio ultima fiat

% A544-P0072-D04
\[
k=\lambda l+1
\]

% A544-P0072-D05
Quum manifesto habeatur

% A544-P0072-D06
\[
\begin{aligned}
\gbr{\frac ab}&=\gbr{\beta+\frac cb}=\beta+\gbr{\frac cb}\\
\gbr{\frac{2a}b}&=\gbr{2\beta+\frac{2c}b}=2\beta+\gbr{\frac{2c}b}\\
\gbr{\frac{3a}b}&=\gbr{3\beta+\frac{3c}b}=3\beta+\gbr{\frac{3c}b}
\end{aligned}
\]

% A544-P0072-D07
etc.. erit

% A544-P0072-D08
\[
\varphi(b,a)=\varphi(b,c)+\tfrac12\beta(b'b'+b')
\]

% A544-P0072-D09
et proin

% A544-P0072-D10
\[
\varphi(a,b)=a'b'-\tfrac12\beta(b'b'+b')-\varphi(b,c)
\]

% A544-P0072-D11
Per similia ratiocinia fit, si statuimus $[\frac12c]=c',\ [\frac12d]=d'\ [\frac12e]=e'$ etc.,

% A544-P0072-D12
\[
\begin{aligned}
\varphi(b,c)&=b'c'-\tfrac12\gamma(c'c'+c')-\varphi(c,d)\\
\varphi(c,d)&=c'd'-\tfrac12\delta(d'd'+d')-\varphi(d,e)\mathord{.}\\
\varphi(d,e)&=d'e'-\tfrac12\varepsilon(e'e'+e')-\varphi(e,f)
\end{aligned}
\]

% A544-P0072-D13
etc. usque ad

% A544-P0072-D14
\[
\varphi(k,l)=k'l'-\tfrac12\lambda(l'l'+l')-\varphi(l,1)
\]

% A544-P0072-D15
Hinc, quoniam manifesto est $\varphi(l,1)=0$, colligimus formulam

% A544-P0072-D16
\[
\begin{aligned}
\varphi(a,b)&=a'b'-b'c'+c'd'-d'e'+\text{etc.}\pm k'l'\\
&\quad-\tfrac12\beta(b'b'+b')+\tfrac12\gamma(c'c'+c')-\tfrac12\delta(d'd'+d')\\
&\quad+\tfrac12\varepsilon(e'e'+e')-\text{etc.}\mp\tfrac12\lambda(l'l'+l')
\end{aligned}
\]

% A544-P0072-D17
\Article{4}

% A544-P0072-D18
Facile iam ex iis, quae in demonstratione tertia exposita sunt, colligitur, relationem numeri $b$ ad $a$, quoties $a$ sit numerus primus, sponte cognosci e va%
\SourceLeaf{73}%
% A544-P0073-D02
lore aggregati $\varphi(a,2b)$. Scilicet prout hoc aggregatum est numerus par vel impar, erit $b$ residuum quadraticum ipsius $a$ vel non-residuum. Ad eundem vero finem ipsum quoque aggregatum $\varphi(a,b)$ adhiberi poterit, ea tamen restrictione, ut casus ubi $b$ impar est ab eo ubi par est distinguatur. Scilicet

% A544-P0073-D03
I.\quad Quoties $b$ est impar, erit $b$ residuum vel non-residuum quadraticum ipsius $a$, prout $\varphi(a,b)$ par est vel impar.

% A544-P0073-D04
II.\quad Quoties $b$ est par, eadem regula valebit, si insuper $a$ est vel formae $8n+1$ vel formae $8n+7$; si vero pro valore pari ipsius $b$ modulus $a$ est vel formae $8n+3$ vel formae $8n+5$, regula opposita applicanda erit, puta, $b$ erit residuum quadraticum ipsius $a$, si $\varphi(a,b)$ est impar, non-residuum vero, si $\varphi(a,b)$ est par.

% A544-P0073-D05
Haec omnia ex art. 4 demonstrationis tertiae facillime derivantur.

% A544-P0073-D06
\Article{5}

% A544-P0073-D07
\textit{Exemplum.} Si quaeritur relatio numeri $103$ ad numerum primum $379$, habemus, ad eruendum aggregatum $\varphi(379,103)$,

% A544-P0073-D08
\[
\begin{array}{r|r|l}
a=379&a'=189&\\
b=103&b'=51&\beta=3\\
c=70&c'=35&\gamma=1\\
d=33&d'=16&\delta=2\\
e=4&e'=2&\varepsilon=8
\end{array}
\]

% A544-P0073-D09
hinc

% A544-P0073-D10
\[
\begin{aligned}
\varphi(379,103)&=9639-1785+560-32-3978+630-272+24\\
&=4786
\end{aligned}
\]

% A544-P0073-D11
unde $103$ erit residuum quadraticum numeri $379$. Si ad eundem finem aggregatum $(379,206)$ adhibere malumus, habemus hocce paradigma:

% A544-P0073-D12
\[
\begin{array}{r|r|l}
379&189&\\
206&103&1\\
173&86&1\\
33&16&5\\
8&4&4
\end{array}
\]

% A544-P0073-D13
unde deducimus

\SourceLeaf{74}%
% A544-P0074-D02
\[
\begin{aligned}
\varphi(379,206)&=19467-8858+1376-64-5356+3741-680+40\\
&=9666
\end{aligned}
\]

% A544-P0074-D03
quapropter $103$ est residuum quadraticum numeri $379$.

% A544-P0074-D04
\Article{6}

% A544-P0074-D05
Quum ad decidendam relationem numeri $b$ ad $a$ non opus sit, singulas partes aggregati $\varphi(a,b)$ computare, sed sufficiat novisse, quot inter eas sint impares, regula nostra ita quoque exhiberi potest:

% A544-P0074-D06
Fiat ut supra $a=\beta b+c,\ b=\gamma c+d,\ c=\delta d+e$ etc., donec in serie numerorum $a,\ b,\ c,\ d,\ e$ etc. ad unitatem perventum sit. Statuatur $[\frac12a]=a',\ [\frac12b]=b',\ [\frac12c]=c'$ etc., sitque $\mu$ multitudo numerorum imparium in serie $a',b',c'$ etc. eorum, quos immediate sequitur impar; sit porro $\nu$ multitudo numerorum imparium in serie $\beta,\gamma,\delta$ etc. eorum, quibus in serie $b',c',d'$ etc. resp. respondent numerus formae $4n+1$ vel formae $4n+2$. His ita factis, erit $b$ residuum quadraticum vel non-residuum ipsius $a$, prout $\mu+\nu$ est par vel impar, unico casu excepto, ubi simul est $b$ par atque $a$ vel formae $8n+3$ vel $8n+5$, pro quo regula opposita valet.

% A544-P0074-D07
In exemplo nostro series $a',b',c',d',e'$ duas successiones imparium sistit, unde $\mu=2$; in serie $\beta',\gamma',\delta',\varepsilon'$ duo quidem impares adsunt, sed quibus in serie $b',c',d',e'$ respondent numeri formae $4n+3$, unde $\nu=0$. Fit itaque $\mu+\nu$ par, adeoque $103$ residuum quadraticum numeri $379$.

% A544-P0074-D08
\BookRule{42mm}

